\documentclass[11pt,a4paper]{article}
\usepackage[utf8]{inputenc}
\usepackage[T1]{fontenc}
\usepackage[french]{babel}
\usepackage{lmodern}
\usepackage{geometry}
\usepackage{xcolor}
\usepackage{hyperref}
\usepackage{listings}
\usepackage{booktabs}
\usepackage{enumitem}
\geometry{margin=2.2cm}
\hypersetup{colorlinks=true,urlcolor=blue,linkcolor=black}
\definecolor{codebg}{RGB}{245,247,250}
\lstset{
  basicstyle=\ttfamily\small,
  backgroundcolor=\color{codebg},
  frame=single,
  breaklines=true,
  columns=fullflexible,
  showstringspaces=false
}

\title{Travaux Pratiques : Infrastructure as Code\\et pipeline CI/CD}
\author{Mehdi\\Dépôt : \url{https://github.com/Mehdi4070/TP3_DevOps}}
\date{16 novembre 2026}

\begin{document}
\maketitle
\tableofcontents
\newpage

\section{Objectif}

L'objectif du TP est de décrire une infrastructure locale sous forme de code avec
Terraform, puis d'automatiser son cycle de déploiement dans GitHub Actions.
L'infrastructure obtenue contient :

\begin{itemize}[noitemsep]
  \item un conteneur PostgreSQL construit à partir de \texttt{postgres:latest} ;
  \item une application Web Nginx construite depuis \texttt{Dockerfile\_app} ;
  \item le port PostgreSQL interne 5432 publié sur le port 5432 ;
  \item le port HTTP interne 80 publié sur le port 8080.
\end{itemize}

\section{Structure et choix techniques}

\begin{tabular}{@{}ll@{}}
\toprule
Fichier & Rôle \\
\midrule
\texttt{main.tf} & Provider Docker, images et conteneurs \\
\texttt{variables.tf} & Identifiants PostgreSQL et ports configurables \\
\texttt{outputs.tf} & URL, endpoint et noms des ressources \\
\texttt{Dockerfile\_app} & Image Alpine légère avec Nginx \\
\texttt{nginx.default.conf} & Virtual host servant la page de validation \\
\texttt{.github/workflows/main.yml} & Pipeline Terraform automatique \\
\texttt{.gitignore} & Exclusion de l'état et des fichiers sensibles \\
\bottomrule
\end{tabular}

Le provider utilisé est \texttt{kreuzwerker/docker} en version compatible
\texttt{\textasciitilde> 3.0}. Le mot de passe est déclaré comme variable sensible afin
qu'il soit masqué dans les sorties Terraform. L'application dépend explicitement du
conteneur PostgreSQL avec \texttt{depends\_on}.

\section{Conteneurisation de l'application}

Le fichier \texttt{Dockerfile\_app} utilise \texttt{alpine:latest}, installe Nginx sans
conserver le cache APK, crée la page demandée et démarre Nginx au premier plan.

\begin{lstlisting}
FROM alpine:latest
RUN apk add --no-cache nginx
RUN mkdir -p /var/www/localhost/htdocs /run/nginx && \
    echo "<h1>Application Deployed via Terraform IaC!</h1>" \
    > /var/www/localhost/htdocs/index.html
EXPOSE 80
CMD ["nginx", "-g", "daemon off;"]
\end{lstlisting}

\section{Exécution du cycle de vie du déploiement}

Les validations ci-dessous ont été exécutées le 16 novembre 2026 avec Terraform
1.17.2 et le provider Docker 3.9.0. Les ports 8080 et 5432 étaient déjà occupés par
d'autres services sur le poste de test. Pour ne pas interrompre ces services, le test
local a utilisé temporairement 18080 et 15432 via \texttt{-var}. Les valeurs par défaut
du projet restent bien 8080 et 5432 ; elles sont utilisées et validées dans GitHub
Actions.

\subsection{Initialisation}

\begin{lstlisting}
terraform init -input=false
\end{lstlisting}

Résultat : backend initialisé et provider
\texttt{kreuzwerker/docker v3.9.0} installé avec succès. La configuration a ensuite été
contrôlée avec \texttt{terraform fmt -check} et \texttt{terraform validate} ; le résultat
était \texttt{Success! The configuration is valid.}

\subsection{Planification}

\begin{lstlisting}
terraform plan -out=tfplan \
  -var='postgres_external_port=15432' \
  -var='app_external_port=18080'
\end{lstlisting}

Résultat : \texttt{Plan: 4 to add, 0 to change, 0 to destroy}. Les quatre ressources
prévues étaient les deux images et les deux conteneurs Docker.

\subsection{Application}

\begin{lstlisting}
terraform apply -auto-approve tfplan
\end{lstlisting}

Résultat : les conteneurs \texttt{tp-iac-postgres} et \texttt{tp-iac-webapp} ont été
créés et démarrés. Les sorties confirmaient notamment :

\begin{lstlisting}
postgres_container_name = "tp-iac-postgres"
postgres_endpoint       = "localhost:15432"
web_app_container_name  = "tp-iac-webapp"
web_app_url             = "http://localhost:18080"
\end{lstlisting}

Un second \texttt{terraform plan} a retourné \texttt{No changes}, ce qui démontre
l'idempotence et la concordance entre l'état réel et la configuration.

\subsection{Validation des services}

La requête HTTP a retourné \texttt{HTTP/1.1 200 OK} ainsi que le contenu exact demandé :

\begin{lstlisting}
<h1>Application Deployed via Terraform IaC!</h1>
\end{lstlisting}

La base a été contrôlée depuis son conteneur :

\begin{lstlisting}
docker exec tp-iac-postgres \
  pg_isready -U tp_user -d tp_database
/var/run/postgresql:5432 - accepting connections
\end{lstlisting}

\subsection{Destruction}

\begin{lstlisting}
terraform destroy -auto-approve \
  -var='postgres_external_port=15432' \
  -var='app_external_port=18080'
\end{lstlisting}

Résultat : \texttt{Destroy complete! Resources: 4 destroyed.} La liste des ressources
de l'état et la recherche des conteneurs du TP étaient toutes deux vides après la
commande.

\section{Pipeline CI/CD}

Le workflow \texttt{.github/workflows/main.yml} se déclenche automatiquement lors de
chaque \texttt{push} vers \texttt{main}. Il exécute dans l'ordre : checkout, installation
de Terraform, contrôle du format, Init, Validate, Plan, Apply, validation HTTP,
validation PostgreSQL et affichage des outputs.

Le plan est enregistré dans \texttt{tfplan} puis fourni à Apply. De plus, Apply possède
la condition explicite \texttt{if: success()}, donc il ne peut être exécuté que si toutes
les étapes précédentes, en particulier Plan, ont réussi. Le workflow teste aussi le
texte de la page sur \url{http://localhost:8080/} et l'état de PostgreSQL.

La variable \texttt{postgres\_password} a fait l'objet d'une modification dédiée dans
l'historique initial du projet et le pipeline associé avait réussi. La version finale
est vérifiable dans l'onglet Actions du dépôt :
\url{https://github.com/Mehdi4070/TP3_DevOps/actions}.

\section{Réponses aux questions d'approfondissement}

\subsection{Impact de \texttt{terraform destroy} sur l'état}

\texttt{terraform.tfstate} est l'inventaire de correspondance entre les adresses des
ressources Terraform et les objets réellement gérés. Lors d'un Destroy, Terraform lit
cet état pour identifier les objets à supprimer, demande leur destruction au provider,
puis retire de l'état les ressources supprimées. Le fichier peut encore exister après
l'opération, mais il ne contient plus les ressources détruites. L'état permet donc la
réconciliation : Terraform compare la configuration souhaitée, l'état connu et les
informations relues depuis l'infrastructure afin de calculer les actions nécessaires.

\subsection{Immuabilité de l'infrastructure}

L'immuabilité signifie qu'une ressource déployée n'est pas modifiée manuellement et
progressivement. Lorsqu'un changement incompatible est nécessaire, on produit une
nouvelle version décrite par le code et on remplace l'ancienne ressource. Cette approche
réduit la dérive de configuration, rend les déploiements reproductibles et facilite le
retour à une version précédente. Terraform décide selon les capacités du provider si un
champ peut être modifié sur place ou s'il impose un remplacement.

\subsection{Pourquoi Plan est une gate essentielle}

Plan affiche à l'avance les créations, modifications, remplacements et destructions. Il
permet de détecter une erreur de variable, une suppression inattendue, une dérive ou un
impact trop important avant tout changement réel. En production, le plan peut être
archivé et soumis à validation humaine ou à des règles automatiques. Appliquer le plan
enregistré garantit ensuite que l'action validée est bien celle qui est exécutée.

\subsection{Création d'un réseau Docker}

La ressource à ajouter est \texttt{docker\_network}. Il faudrait ensuite connecter les
deux conteneurs à ce réseau, par exemple avec leur bloc \texttt{networks\_advanced}.
L'application pourrait alors joindre PostgreSQL par son nom ou son alias sur le réseau
privé, sans dépendre du port publié sur l'hôte.

\section{Conclusion}

Le TP met en œuvre le cycle complet Init, Plan, Apply, Validate et Destroy. La
configuration est valide, reproductible et idempotente. L'application et PostgreSQL ont
été testés réellement, et le même processus est automatisé sur chaque push vers
\texttt{main} avec une gate explicite avant Apply.

\end{document}
